\SourceLeaf{173}
\Article{36}
Ein weiterer Punkt unserer Theorie, an dem Zweifel bestehen können, ist die Annahme, dass die Ursachen der erdmagnetischen Kraft ausschließlich im Inneren der Erde liegen.

Sollten die unmittelbaren Ursachen ganz oder teilweise außerhalb der Erde zu suchen sein, so können wir. sofern wir haltlose Spekulationen ausschließen und uns nur an wissenschaftlich Bekanntes halten wollen, nur an galvanische Ströme denken. Die Luft der Atmosphäre leitet solche Ströme ebenso wenig wie der leere Raum. Unser Wissen lässt uns also im Stich, wenn wir in den oberen Regionen nach einem Träger für galvanische Ströme suchen. Doch die rätselhaften Erscheinungen des Nordlichts, bei denen allem Anschein nach Elektrizität in Bewegung eine Hauptrolle spielt, verbieten es uns, die Möglichkeit solcher Ströme allein wegen dieses Nichtwissens rundweg zu leugnen. Jedenfalls bleibt es interessant zu untersuchen, welche magnetische Wirkung diese Ströme auf der Erdoberfläche hervorbringen würden.

\Article{37}
Nehmen wir also an, dass in einem Raum $S$, der die Erde wie ein Gewölbe oder eine Schale umschließt, beständige galvanische Ströme fließen. Den gesamten von $S$ eingeschlossenen Raum bezeichnen wir mit $S'$ und den gesamten Außenraum, der $S$ und $S'$ umschließt, mit $S''$. Wie diese galvanischen Ströme auch angeordnet sein mögen, sie lassen sich stets durch eine fiktive Verteilung magnetischer Flüssigkeiten innerhalb des Raumes $S$ selbst ersetzen, die im gesamten übrigen Raum $S'$ und $S''$ genau dieselbe magnetische Wirkung ausübt wie die Ströme. Dieser wichtige, bereits in Artikel 3 erwähnte Satz beruht erstens darauf, dass sich die Ströme in unendlich viele elementare Ströme zerlegen lassen (das heißt in solche, die als linear betrachtet werden dürfen). Zweitens beruht er auf dem bekannten, meines Wissens zuerst von \textsc{Ampère} bewiesenen Theorem, nach dem jeder lineare Strom, der eine beliebige Fläche begrenzt, durch eine Verteilung der magnetischen Flüssigkeiten auf beiden Seiten dieser Fläche in unmessbar kleinen Abständen von ihr mit der genannten Wirkung ersetzt werden kann. Drittens beruht er auf der offensichtlichen Möglichkeit, zu jeder innerhalb von $S$ liegenden geschlossenen Linie eine von ihr begrenzte Fläche anzugeben, die ebenfalls ganz innerhalb von $S$ liegt.

\SourceLeaf{174}
Bezeichnet man nun mit $-v$ die Summe aller Quotienten, die man erhält, wenn man sämtliche Elemente dieses fiktiven magnetischen Fluids jeweils durch ihren Abstand von einem unbestimmten Punkt $O$ in $S'$ oder $S''$ dividiert und dabei die Elemente des südlichen Fluids selbstverständlich als negativ ansetzt, so drücken die partiellen Ableitungen von $v$ (genau wie die von $V$ in unserer obigen Theorie) die Komponenten der magnetischen Kraft aus, die die galvanischen Ströme in $O$ hervorbringen.

\Article{38}
Die ausführliche Entwicklung der Theorie, aus der der im vorigen Artikel verwendete Satz stammt, muss einer anderen Gelegenheit vorbehalten bleiben. Ein wichtiger Punkt dieser Theorie verdient aber noch, hier erwähnt zu werden. Konstruiert man zwei \textit{verschiedene} Flächen $F$, $F'$, die beide von demselben linearen Strom $G$ begrenzt werden, und betrachtet man der Kürze halber nur den einfachsten Fall, in dem die Flächen außer ihrer gemeinsamen Randlinie keinen weiteren Punkt gemein haben, so schließen sie einen dreidimensionalen Raum ein. Liegt $O$ außerhalb dieses Raumes, so erhält man für denjenigen Anteil von $v$, der sich auf $G$ bezieht, \textit{ein und denselben} Wert, ob man die magnetischen Fluide auf $F$ oder auf $F'$ verteilt. Er ist gleich dem Produkt aus der Stärke des galvanischen Stroms $G$ (gemessen in einer geeigneten Einheit) und dem Raumwinkel, dessen Scheitel in $O$ liegt und der von den Geraden von $O$ zu den Punkten von $G$ eingeschlossen wird. Das ist dasselbe wie das Produkt aus der Stromstärke und dem Teil der Kugelfläche mit Radius 1 um $O$, der die gemeinsame Projektion von $F$ und $F'$ ist. Liegt dagegen $O$ innerhalb des von $F$ und $F'$ eingeschlossenen Raumes, so sind die beiden Werte des betreffenden Anteils von $v$ verschieden, je nachdem, ob man die magnetischen Flüssigkeiten auf $F$ oder auf $F'$ verteilt. Ihnen entsprechen nämlich verschiedene Teile der genannten Kugelfläche, die sich zur ganzen Kugelfläche ergänzen. Da die Richtung des galvanischen Stroms gegenüber $F$ und $F'$ entgegengesetzt ist, muss man aber in beiden Fällen der Stromstärke bei der Multiplikation mit den Teilen der Kugelfläche entgegengesetzte Vorzeichen geben. Daraus folgt, dass die algebraische Differenz zwischen den beiden Werten
\SourceLeaf{175}
des betreffenden Anteils von $v$ gleich der Stromstärke multipliziert mit der gesamten Kugelfläche, also mit $4\pi$, ist.

Daraus folgt leicht, dass, wenn $O$ in $S''$ liegt, der Wert von $v$ von der Wahl der Verbindungsflächen ganz unabhängig bleibt. Liegt $O$ dagegen in $S'$, so hängt zwar der absolute Wert von $v$ von dieser Wahl ab, die Differentiale von $v$ bleiben aber davon unabhängig.

Das hier angesprochene, höchst fruchtbare Theorem bedarf in seiner Allgemeinheit noch weiterer Erläuterungen. Ihm zufolge hat für die magnetische Wirkung eines linearen galvanischen Stroms das Produkt aus seiner Stärke und dem Teil der Kugelfläche, den die Projektion der Stromlinie von $O$ aus begrenzt, dieselbe Bedeutung, die für Anziehungs- oder Abstoßungskräfte die durch den Abstand von $O$ dividierten Massenteile haben. Diese Erläuterungen müssen einer ausführlichen Behandlung des Gegenstandes vorbehalten bleiben.
